\documentclass[11pt]{article}
\usepackage[T1]{fontenc}
\usepackage{lmodern,microtype,amsmath,amssymb,amsthm,booktabs,tabularx,array}
\usepackage[a4paper,top=25mm,bottom=25mm,left=27mm,right=27mm]{geometry}
\usepackage{xcolor,fancyhdr,enumitem}
\definecolor{ink}{HTML}{173E50}
\usepackage[colorlinks=true,linkcolor=ink,citecolor=ink,urlcolor=ink,pdftitle={Sharp Randic bounds and the structure of extremal matchings}]{hyperref}
\newtheorem{theorem}{Theorem}[section]
\newtheorem{lemma}[theorem]{Lemma}
\newtheorem{proposition}[theorem]{Proposition}
\newtheorem{corollary}[theorem]{Corollary}
\theoremstyle{remark}\newtheorem{remark}[theorem]{Remark}
\setcounter{tocdepth}{1}
\setlist{nosep}
\setlength{\emergencystretch}{2em}
\pagestyle{fancy}\fancyhf{}\lhead{\small\textsc{Randi\'c index and matchings}}\rhead{\small Research manuscript}\cfoot{\thepage}\setlength{\headheight}{14pt}

\begin{document}
\begin{center}
{\small\color{ink}\textsc{Extremal graph theory \enspace / \enspace research manuscript}}\\[10pt]
{\LARGE\bfseries Sharp Randi\'c bounds and the structure\\[3pt] of extremal matchings}\\[9pt]
{\normalsize With companion results on fractional curvature and $p$-adic factorial sums}\\[8pt]
{\small Prepared for Keio \enspace$\cdot$\enspace 16 November 2026}
\end{center}
\begin{abstract}
For every finite simple graph $G$, its Randi\'c index is at least the two-thirds power of its fractional matching number. The constant one is sharp. For ordinary matching number $k$ tending to infinity the infimum of $R(G)/k^{2/3}$ tends to $3\,2^{-1/3}$. This proves Conjecture~9 of Akbari et al.\ and determines the limiting coefficient asked for in their subsequent remark. An averaging argument selects complementary degree ranges and separates the weight of matched edges from that of the cut. Near equality forces a complete bipartite graph with part sizes asymptotic to $2k$ and $(2k)^{1/3}$ up to $o(k^{4/3})$ edge edits. Two independent applications of corpus-inspired estimates give exact finite curvature interpolation under strictly decreasing polynomial forcing, and the full root count for John C's alternating factorial Mahler series, including an exact dyadic valuation formula.
\end{abstract}
\tableofcontents
\section*{Questions and antecedents}
The principal question is how small the Randi\'c index can be when only the matching number is fixed. Akbari et al.\ \cite{AGGHT} constructed generalized windmills and proved a lower bound for graphs with a nearly perfect matching. Their Conjecture~9 quantifies over all finite simple graphs. Their remark asks how the optimal coefficient depends on the matching number. These are distinct questions: an extremizer at $k=1$ determines the best universal coefficient but gives no information about the limit. The two coefficients below are $1$ and $3\,2^{-1/3}$, respectively.
\par
In VMRC [183] weight is distributed over incidences in a line graph. In [189], Paper~01 sums inequalities across thresholds on vertex degree to control a radius. Here a fractional matching supplies edge weights, and both endpoint degrees determine which thresholds to use. The auxiliary problem is to partition the edges of a near-minimizer so that the two leading contributions can be bounded without counting any edge twice. The answer produces the limiting constant and an edit-distance stability statement. The companion sections can be read independently of the graph proofs.
\clearpage
\section{A sharp Randi\'c bound for fractional matchings}\label{graph:section}

For a finite simple graph $G=(V,E)$ with degrees $d_v$, set
\[
 R(G)=\sum_{uv\in E}(d_ud_v)^{-1/2}.
\]
An ordinary matching is a family of vertex-disjoint edges, and $\nu(G)$
denotes its maximum size. Akbari, Ghasemi Nezhad, Ghazizadeh, Haslegrave
and Tohidi ask in Conjecture~9 of their August 2024 preprint
\cite{AGGHT} whether $R(G)\ge c\,\nu(G)^{2/3}$ for some absolute $c>0$.
They propose $c=1$ immediately after the conjecture. They obtain the
correct exponent for graphs with a nearly perfect matching, where at
most one vertex is unmatched. We remove this restriction and determine
the sharp universal constant.

A fractional matching assigns each edge $e$ a weight $0\le x_e\le1$,
subject to $\sum_{e\ni v}x_e\le1$ at every vertex. Its mass is
$k=\sum_e x_e$, and the largest possible mass is denoted by
$\nu_f(G)$. The indicator of an ordinary matching is feasible.

\begin{theorem}\label{graph:main}
Every finite simple graph satisfies
\[
             R(G)\ \ge\ \nu_f(G)^{2/3}\ \ge\ \nu(G)^{2/3}.
\]
In either lower bound on $R(G)$, for graphs with at least one edge,
equality holds precisely when the deletion of all isolated vertices
leaves $K_2$. In particular the coefficient $1$ is optimal.
\end{theorem}

\subsection{An auxiliary problem on weighted edge sets}

Given a fractional matching of mass $k$ in a graph with $m$ edges,
how much Randi\'c weight is forced separately by its weights $x_e$ and
the complementary weights $1-x_e$? The vertex constraints bound a
first degree moment for the former. Simplicity bounds a degree-product
moment for the latter. Unlike a bound on the maximum vertex degree,
these two estimates are independent of the number of unmatched vertices.
The use of local edge weights is related to the incidence allocations
in corpus entry~183 \cite{VMRC}. Here the allocation takes place in the
original graph and retains both portions of every edge.

\begin{lemma}\label{graph:split}
If $m>0$ and $x$ is a fractional matching of mass $k$, then
\begin{equation}\label{graph:bound}
 R(G)\ge \frac{k^2}{m}+\frac{(m-k)^{3/2}}{\sqrt{2}\,m}.
\end{equation}
\end{lemma}
\begin{proof}
Write $w_{uv}=(d_ud_v)^{-1/2}$. The endpoint constraints and the
handshaking identity give
\[
 \sum_{uv\in E}x_{uv}(d_u+d_v)
 =\sum_{v\in V} d_v\sum_{e\ni v}x_e\le 2m .
\]
If $k>0$, the arithmetic--geometric mean and Cauchy--Schwarz inequalities give
\[
 \sum_{uv\in E}x_{uv}w_{uv}
 \ge \sum_{uv\in E}\frac{2x_{uv}}{d_u+d_v}
 \ge \frac{2k^2}{\sum_{uv\in E}x_{uv}(d_u+d_v)}
 \ge \frac{k^2}{m}.
\]
The same lower bound holds trivially for $k=0$. Set $y_e=1-x_e$,
of total mass $b=m-k$. Simplicity gives
\[
 \sum_{uv\in E}y_{uv}d_ud_v\le
 \sum_{\{u,v\}\subseteq V}d_ud_v
 =\frac{(2m)^2-\sum_v d_v^2}{2}\le2m^2 .
\]
For $b>0$, Jensen's inequality for $t\mapsto t^{-1/2}$ with
probabilities $y_e/b$ gives
\[
 \sum_{uv\in E}y_{uv}w_{uv}\ge
 \frac{b^{3/2}}{\sqrt{\sum_{uv\in E}y_{uv}d_ud_v}}
 \ge \frac{b^{3/2}}{\sqrt{2}\,m}.
\]
For $b=0$ this contribution is zero. Adding proves \eqref{graph:bound}.
\end{proof}

\begin{proof}[Proof of Theorem~\ref{graph:main}]
First let $G$ be connected with at least one edge. Fix a maximum
fractional matching, of mass $k\ge1$.
If $k\ge3$, put $t=m/k^{4/3}$ and $u=k^{-1/3}$. Dividing
\eqref{graph:bound} by $k^{2/3}$ yields
\[
 \frac{R(G)}{k^{2/3}}\ge\frac{1+(t-u)^{3/2}/\sqrt{2}}{t}.
\]
If $t\le1$, the right side exceeds $1$ (at $t=1$ use $u<1$).
For $t>1$ set $z=t-1$ and $c=1-u$. The needed strict inequality
is $(z+c)^3>2z^2$. Applying arithmetic--geometric mean to
$z/2,z/2,c$ gives $(z+c)^3\ge 27cz^2/4$.
As $k\ge3$, $c\ge1-3^{-1/3}>8/27$; the last comparison follows
from $27^3<3\cdot19^3$.

If $1\le k\le3$, summing the vertex constraints gives
$2k\le |V|$. The connected-graph lower bound of Bollob\'as and
Erd\H{o}s \cite{BE} gives $R(G)\ge\sqrt{|V|-1}\ge\sqrt{2k-1}$.
The function $s\mapsto 2s-1-s^{4/3}$ is concave on $[1,3]$,
vanishes at $1$, and is positive at $3$ since $3^4<5^3$.
It is strictly positive for $1<s\le3$. Thus equality can hold only
at $k=1$, where $R(G)=1$ forces $|V|=2$ by the same order bound.
Conversely $K_2$ realizes equality.

Both $R$ and $\nu_f$ add over connected components.
For nonnegative $a_i$, $\sum_i a_i^{2/3}\ge(\sum_i a_i)^{2/3}$,
strictly so if two of them are positive. This gives the claimed bound
and equality classification for arbitrary graphs. An edgeless graph
has both parameters zero. Finally $\nu\le\nu_f$ gives the ordinary
matching result with the same equality condition.
\end{proof}


For an odd cycle the fractional matching can be larger than every
ordinary matching. On a triangle, putting $x_e=1/2$ on each edge
gives $\nu_f=3/2$ and $\nu=1$. The fractionally weighted statement is
therefore a genuine extension. Liu et al.\ \cite{Liu} give an upper
bound in terms of order and matching number for K\H{o}nig--Egerv\'ary
graphs. That result addresses a different extremal question.
\subsection{Exact limiting constant for ordinary matchings}

Write $c_k=\inf_{\nu(G)=k}R(G)/k^{2/3}$, as in the question
following Conjecture~9 of \cite{AGGHT}. We determine its limit.
The proof uses the fact that, in a graph near the minimum, all but
a lower-order number of edges join vertices of complementary degree
scales. We retain the matching contribution even when its edges belong
to that lower-order set.

\begin{theorem}\label{graph:limit}
For finite simple graphs with ordinary matching number $k$,
\[
              \lim_{k\to\infty}c_k=\frac32\,2^{2/3}.
\]
More generally, for any sequence of fractional matchings with
masses $k\to\infty$ in finite simple graphs,
$\liminf R(G)/k^{2/3}\ge (3/2)2^{2/3}$.
\end{theorem}
\begin{proof}
For the lower bound it suffices to consider sequences satisfying
$R(G)\le Ck^{2/3}$ for some fixed $C$. The all-edge Jensen estimate
in the proof of Lemma~\ref{graph:split}, with every $y_e=1$, gives
$R(G)\ge\sqrt{m/2}$. Therefore $m=O(k^{4/3})$. We may ignore isolated
vertices. Set $K=k^{4/3}$ and $q_e=d_ud_v$ for $e=uv$.
For $z\ge0$ the two moment bounds give
\[
 \#\{q_e<Ke^{-z}\}\le R(G)\sqrt K e^{-z/2},\qquad
 \#\{q_e>Ke^{z}\}\le 2m^2 K^{-1} e^{-z}.
\]
Integrating the two inequalities in $z$ gives the estimate
\begin{equation}\label{eq:logcost}
       \sum_{e\in E}\big|\log(q_e/K)\big|=O(K).
\end{equation}
The implied constant depends only on $C$. This controls the total
cost of pairing the logarithms of the endpoint degrees about
$\tfrac12\log K$.

Choose $s$ uniformly on $I=[(2/5)\log k,(1/2)\log k]$ and set
\[
 A=\{v:\log d_v<s\},\qquad H=\{v:\log d_v>\log K-s\}.
\]
The sets are disjoint. For an oriented edge $(u,v)$, membership of
$u$ in $A$ is the event $s>\log d_u$; membership of $v$ in $H$ is the
event $s>\log K-\log d_v$. These indicators disagree with probability
at most $|\log(q_{uv}/K)|/|I|$. Write $D_X=\sum_{v\in X}d_v$ and
$\ell=e(A,H)$. Summing disagreements over both orientations of every
edge counts exactly $D_A+D_H-2\ell$. Thus \eqref{eq:logcost} supplies
a threshold such that
\begin{equation}\label{eq:cuterror}
  0\le D_A-\ell,\ D_H-\ell,\qquad
  (D_A-\ell)+(D_H-\ell)=O(K/\log k).
\end{equation}
This bound selects the cut from the graph itself and controls the error
even when the matching leaves arbitrarily many vertices uncovered.

Outside $A$ there are at most $2m/k^{2/5}=O(k^{14/15})=o(k)$ vertices.
Thus the fractional matching mass $k_A$ on $A$-internal edges is
$k-o(k)$ by the vertex constraints. Its Randi\'c contribution is
at least $2k_A^2/D_A$ by the first estimate in
Lemma~\ref{graph:split}, restricted to $A$. As $R(G)=O(k^{2/3})$ this
forces $D_A=\Omega(K)$ and hence $\ell=\Omega(K)$.
All crossing edges between $A$ and $H$ are disjoint from the
$A$-internal edge set. Their degree-product sum is at most $D_A D_H$
by simplicity, so Jensen gives a contribution of at least
$\ell^{3/2}/\sqrt{D_A D_H}$. Consequently
\[
 R(G)\ \ge\ (1-o(1))
 \left(\frac{2k^2}{\ell}+\sqrt{\ell}\right)
 \ \ge\ (1-o(1))\frac32(2k)^{2/3}.
\]
The last bound follows by arithmetic--geometric mean on
$2k^2/\ell,\sqrt{\ell}/2,\sqrt{\ell}/2$.

In fact, the relative errors from \eqref{eq:cuterror} are
$O(1/\log k)$ since $\ell=\Omega(K)$, and $k-k_A=O(k^{14/15})$.
Taking, for example, $C=3$ and observing that $3>(3/2)2^{2/3}$ gives
an absolute implied constant in the lower estimate
\[
 R(G)\ge \tfrac32(2k)^{2/3}-O(k^{2/3}/\log k)\qquad(k\to\infty).
\]

For the opposite inequality use the generalized windmills of
\cite[Proposition~7]{AGGHT}. Here is the construction with its
normalization. Take two independent parts of sizes $2(k-r)$ and $2r+1$,
join them completely, then add a perfect matching inside the larger
part, where $r=\lfloor(k/4)^{1/3}\rfloor$ and $k$ is large. The order
$2k+1$ bounds the matching number by $k$. Match $2r$ vertices from the
small part into the endpoints of $r$ paired edges of the large part,
then use the $k-2r$ untouched pairs. The attained size is $k$.
The degrees on the large and small parts are $2r+2$ and
$2(k-r)$, respectively, so
\[
 R=\frac{k-r}{2r+2}
 +\frac{2(k-r)(2r+1)}{\sqrt{(2r+2)2(k-r)}}
 =\bigl(\tfrac32+o(1)\bigr)(2k)^{2/3}.
\]
Together with the lower bound this gives the asserted limit.
\end{proof}
\subsection{Structure of asymptotically extremal graphs}

The sharp leading term also determines the large bipartite subgraph in
every sequence approaching equality. The distance here counts edge
additions and deletions on the given vertex set. Extra vertices in the
comparison graph are isolated.

\begin{theorem}\label{graph:stability}
Let $G_k$ be finite simple graphs with $\nu(G_k)=k\to\infty$ and
$R(G_k)=(3/2+o(1))(2k)^{2/3}$. For every choice of maximum matching
there are disjoint sets $X,Y$ such that $X$ is the union of the
endpoint pairs of $k-o(k)$ edges of that matching,
\[
 |X|=(2+o(1))k,\qquad |Y|=(1+o(1))(2k)^{1/3}.
\]
With $P_X$ denoting those matching edges, one has
\[
 \bigl|E(G_k)\mathbin{\triangle}
       \bigl(E(K_{X,Y})\cup P_X\bigr)\bigr|=o(k^{4/3}).
\]
\end{theorem}
\begin{proof}
Put $K=k^{4/3}$ and take $A,H,\ell$ as in the proof of
Theorem~\ref{graph:limit}. In particular $m=O(K)$,
$D_A,D_H\sim\ell$, and $\ell$ is bounded above and below by positive
constant multiples of $K$. Fix a maximum matching. Write $P$ for its
$k'=k-o(k)$ internal edges of $A$ and $U$ for their endpoints. Define
$W_P$ and $W_{AH}$ as the sums of the Randi\'c edge weights over $P$
and over the $A,H$ cut edges, respectively.

The function $2/t+\sqrt{t}$ has its unique minimum over $t>0$ at
$t=4^{2/3}$. The proof of Theorem~\ref{graph:limit}, applied to
$t=\ell/K$, therefore gives $\ell\sim4^{2/3}K$. Both disjoint lower
bounds used there must be asymptotically attained. In formulas,
\begin{equation}\label{graph:saturation}
 W_P\sim 2k^2/\ell,\quad W_{AH}\sim\sqrt{\ell},\quad
 R(G_k)-W_P-W_{AH}=o(\sqrt K).
\end{equation}
Indeed the sum of the lower bounds is $R(G_k)-o(\sqrt K)$
and each of the two bounds is of order $\sqrt K$.
From $W_P\ge2(k')^2/D_U$ and $D_U\le D_A\sim\ell$ we get
$D_U\sim\ell$. Write $\mu=D_U/(2k')$.

Choose an edge $uv$ uniformly from $P$ and put
$s=(d_u+d_v)/(2\mu)$, $r=(d_u-d_v)/(d_u+d_v)$, with either ordering of
the endpoints. Then $\mathbb E s=1$ and
\[
 \mathbb E\!\left[\frac{1}{s\sqrt{1-r^2}}\right]
      =\frac{\mu}{k'}W_P=1+o(1).
\]
Since $\mathbb E(1/s)\ge1$, this implies
$\mathbb E[(s-1)^2/s]=\mathbb E(1/s)-1=o(1)$, and
$\mathbb E|s-1|=o(1)$ by Cauchy--Schwarz.
It also gives $\mathbb E[(1/\sqrt{1-r^2}-1)/s]=o(1)$.
On $1/2\le s\le 2$ the last expression forces $r\to0$ in probability.
The complementary event has probability $o(1)$, hence
$\mathbb E[s|r|]\le\mathbb E|s-1|+\mathbb E|r|=o(1)$.
The two endpoint degrees divided by $\mu$ are $s(1\pm r)$, so their
average absolute deviations from $1$ tend to zero. It follows
that there is a sequence $\epsilon_k\downarrow0$ such that removing
pairs for which either normalized degree differs from $1$ by more
than $\epsilon_k$ discards $o(k)$ pairs and $o(K)$ degree mass.
For example take the square root of the mean sum of the two
deviations as a threshold, enlarged by $1/k$.
Let $X$ be the endpoints of the remaining pairs. Thus $D_X\sim\ell$,
$|X|=(2+o(1))k$, and every vertex in $X$ has degree $(1+o(1))\mu$
uniformly.

Near equality in the cut Jensen bound gives
\[
 \sum_{uv\in E(A,H)}d_ud_v\ge\frac{\ell^3}{W_{AH}^2}
             =(1-o(1))D_A D_H.
\]
Consequently the sum of $d_ud_v$ over missing pairs in $A\times H$
is $o(\ell^2)$. Since each $u\in X$ has degree at least
$(1-o(1))\mu$ and $\mu|X|\sim\ell$, this implies
\[
 \sum_{v\in H} d_v\frac{|X\setminus N(v)|}{|X|}=o(D_H).
\]
For some $\eta_k\downarrow0$, the vertices $v\in H$ missing more than
$\eta_k |X|$ neighbors of $X$ have total degree $o(D_H)$, by Markov's
inequality. Discard them and call the remaining set $Y$. We have
$D_{A\setminus X}=o(\ell)$ and $D_{H\setminus Y}=o(\ell)$, so
$e(X,Y)=(1-o(1))\ell$. Every $v\in Y$ is adjacent to at least
$(1-\eta_k)|X|$ vertices of $X$, whence
\[
 (1-\eta_k)|X||Y|\le e(X,Y)\le |X||Y|.
\]
Thus $|X||Y|\sim\ell$, the number of missing pairs is $o(K)$,
and $|Y|\sim\ell/(2k)=(1+o(1))(2k)^{1/3}$.

Finally, let $N$ be the number of edges outside $P\cup E(A,H)$.
Their degree-product sum is at most $2m^2$, so their Randi\'c weight
is at least $N^{3/2}/(\sqrt{2}m)$ by Jensen. Equation
\eqref{graph:saturation} forces $N=o(K)$ since $m=O(K)$.
There are $o(K)$ edges of $E(A,H)$ outside $E(X,Y)$ and $o(k)$ edges
of $P$ outside $P_X$. Together with the count of missing pairs this
proves the edit-distance assertion.
\end{proof}

\clearpage
\section{Prescribing curvature by decreasing polynomial forcing}
\label{sec:interpolation}

How much curvature information survives the endpoint conditions for a general forcing?
Fix $1<\alpha<2$ and write
\[
 (I^\alpha f)(x)=\frac{1}{\Gamma(\alpha)}
    \int_0^x(x-t)^{\alpha-1} f(t)\,dt,\qquad
 U_f(x)=(I^\alpha f)(x)-x^{\alpha-1}(I^\alpha f)(1).
\]
This is the solution with homogeneous Dirichlet conditions in
\cite[eq.~(7)]{LMS}. VMRC [123] extracts zero information from fractional kernels by estimating their integral terms. Here differentiation exposes a singular contribution that is localized at a prescribed site. The following interpolation problem fixes the
observation sites in advance and allows arbitrary real observations.

\begin{theorem}\label{thm:curvature-interpolation}
For any integer $n\ge1$, points $0<x_1<\cdots<x_n<1$, and vector
$y\in\mathbb R^n$, there is a real polynomial $f$ such that
\[
 f(1)=0,\qquad f'(x)<0\quad(0\le x\le1),\qquad
 U_f''(x_i)=y_i\quad(1\le i\le n).
\]
In particular, $f>0$ on $[0,1)$.
\end{theorem}

\begin{proof}
Set $\mathcal L_i f=U_f''(x_i)$ and $\mathcal L f=(\mathcal L_i f)_{i=1}^n$.
For $f\in C^1[0,1]$ and $0<x\le1$, differentiation after
writing $f(t)=f(0)+\int_0^t f'(s)\,ds$ gives
\begin{align}
 U_f''(x)
 &=\frac{f(0)x^{\alpha-2}}{\Gamma(\alpha-1)}
 +\frac{1}{\Gamma(\alpha-1)}
       \int_0^x(x-t)^{\alpha-2}f'(t)\,dt \notag\\
 &\hspace{15mm}+(\alpha-1)(2-\alpha)x^{\alpha-3}(I^\alpha f)(1).
 \label{eq:interp-eval}
\end{align}
Each $\mathcal L_i$ is continuous on $C^1[0,1]$ since $\alpha-2>-1$.

First choose a polynomial $f_0$ with $f_0(1)=0$, $f_0'<0$ and
$v_i:=\mathcal L_i f_0> y_i$ at every site. To construct it, start with a
smooth nonincreasing function equal to one on $[0,x_n]$ and zero near
one. Its curvature is strictly positive at all the sites by
\eqref{eq:interp-eval}. Adding a sufficiently small positive multiple
of $1-x$ preserves these strict inequalities and makes the derivative
uniformly negative. Approximate the derivative uniformly by a polynomial
and integrate it back from the endpoint value zero. The resulting
polynomial retains all strict signs. A positive scalar multiple gives
the required $f_0$.

Choose $S\in C^\infty(\mathbb R)$ nonincreasing, equal to one on
$(-\infty,-1]$ and zero on $[0,\infty)$. For $\varepsilon>0$ small, set
$g_{j,\varepsilon}(t)=S((t-x_j)/\varepsilon)$. Its entire drop lies
between $x_j-\varepsilon$ and $x_j$. At the $j$th site its contribution
from the integral of $f'$ in \eqref{eq:interp-eval} is
$-d\,\varepsilon^{\alpha-2}$, where
\[
 d=-\frac{1}{\Gamma(\alpha-1)}
        \int_{-1}^{0}(-u)^{\alpha-2}S'(u)\,du>0.
\]
At earlier sites that integral vanishes. At later sites the kernel is
bounded on the support of $g'_{j,\varepsilon}$, whose absolute integral
is one. The other two terms of \eqref{eq:interp-eval} stay bounded
at all sites because $0\le g_{j,\varepsilon}\le1$. Thus
\[
 \lim_{\varepsilon\downarrow0}
   \varepsilon^{2-\alpha}\mathcal L_i g_{j,\varepsilon}
   =-d\,\delta_{ij}.
\]

For each fixed $j,\varepsilon$, approximate
$\sqrt{-g'_{j,\varepsilon}}$ uniformly on $[0,1]$ by a polynomial $p$.
Then $b_{j,\varepsilon}(x)=\int_x^1 p(t)^2\,dt$ is a nonincreasing
polynomial vanishing at one and can approximate $g_{j,\varepsilon}$
arbitrarily closely in $C^1$. Choose it so that
$|\mathcal L_i(b_{j,\varepsilon}-g_{j,\varepsilon})|<1$ for every $i$.
The matrices
\[
 A_\varepsilon=
  \bigl(\varepsilon^{2-\alpha}\mathcal L_i b_{j,\varepsilon}\bigr)_{i,j=1}^n
 \quad\hbox{therefore tend to}\quad -d\,\mathrm{Id}_n.
\]
For sufficiently small $\varepsilon$, they are invertible and
$a=A_\varepsilon^{-1}(y-v)$ has all coordinates positive, since its
limit is $(v-y)/d$. Now take
\[
 f=f_0+\varepsilon^{2-\alpha}\sum_{j=1}^n a_j b_{j,\varepsilon}.
\]
Its derivative is strictly negative, its endpoint value is zero,
and linearity gives $\mathcal L f=v+A_\varepsilon a=y$.
\end{proof}

\begin{corollary}\label{cor:many-inflections}
At every order $1<\alpha<2$ and for every $m\ge1$, a strictly
decreasing positive-on-$[0,1)$ polynomial forcing satisfying all
the hypotheses of \cite[Thm.~5.2]{LMS} yields at least $2m-1$ distinct
sign-changing interior curvature zeros.
\end{corollary}

\begin{proof}
Prescribe alternating values $+1,-1,\ldots,+1,-1$ at $2m$ ordered
interior points. On $(0,1)$ the function $U_f''$ is real analytic:
fractional integration of each monomial gives a constant times a real
power of $x$. On each closed interval between consecutive observation
points it has finitely many zeros and opposite signs at the ends.
At least one zero there changes sign. The polynomial $f$ has bounded
derivative at zero, as required by the weighted endpoint condition.
\end{proof}

The public record \cite{Omega} supplies a specific decreasing-polynomial
counterexample to uniqueness and records a separated-pulse proposal
for unbounded multiplicity. Theorem~\ref{thm:curvature-interpolation}
uses localized downward transitions to prescribe exact values at
predetermined sites. The nearly diagonal matrix in the proof makes
the correction coefficients positive while retaining the monotonicity
of the forcing.


\section{Roots of John C's alternating factorial sum}\label{fac:section}

For a prime $p$ and an integer $k\ge 0$, set
\[
 u_p(k)=(-1)^{\lfloor k/p\rfloor}\frac{(pk)!}{p^k k!},
 \qquad S_p(n)=\sum_{k=1}^n u_p(k),\qquad S_p(0)=0.
\]
In March 2026 John C \cite{John} asked for the number of zeros in $\mathbb Z_p$ of
the following specific Mahler series:
\begin{equation}\label{fac:mahler}
 F_p(x)=\sum_{m\ge1}\binom{x}{m}
 \sum_{j=0}^{m-1}(-1)^{m-1-j}\binom{m-1}{j}u_p(j+1).
\end{equation}
The boxed conjecture in the original question predicts $p$ distinct
zeros for every prime. The surrounding prose predicts a single zero
when $p=2$. The theorem determines which prediction holds and counts
all the other fibers at odd primes as well.

\begin{theorem}\label{fac:main}
Series \eqref{fac:mahler} converges uniformly on $\mathbb Z_p$ and has
exactly $p$ distinct zeros there. For each odd $p$, every $y\in\mathbb Z_p$
has one preimage in every residue class modulo $p$.
For $p=2$, one zero is $0$, and the other is $16$ modulo $32$.
\end{theorem}

\subsection{An auxiliary question about partial sums}
Given an integer sequence $T(n)$, when do its increments at prime-power
steps determine an entire local fiber of its continuous extension?
Write $v_p$ for the additive valuation, with $v_p(0)=+\infty$.

\begin{lemma}\label{fac:fiber}
Fix integers $h,d\ge0$, $0\le r<p^h$, and a unit $b\in\mathbb Z_p^\times$.
Suppose that, for all $j\ge0$ and all $n\ge0$ with $n\equiv r\pmod{p^h}$,
\[
 T(n+p^{h+j})-T(n)\equiv b p^{d+j}\pmod{p^{d+j+1}}.
\]
Then $T$ extends uniquely on $r+p^h\mathbb Z_p$ and maps this ball
bijectively onto $T(r)+p^d\mathbb Z_p$. For $x,x'$ in this ball,
\[
 v_p(T(x)-T(x'))=v_p(x-x')+d-h.
\]
\end{lemma}
\begin{proof}
For nonnegative integers $s<t$, write $t-s=p^j m$ with $p\nmid m$.
Summing $m$ increments starting at $r+p^hs$ gives a difference
congruent to $mbp^{d+j}$ modulo $p^{d+j+1}$.
Thus $(T(r+p^ht)-T(r))/p^d$ is an isometry on the nonnegative integers.
It extends to an isometry on their completion. Its reduction at every
precision is an injection of the finite ring $\mathbb Z/p^k\mathbb Z$
into itself, hence a bijection. The compatible inverses yield
surjectivity on $\mathbb Z_p$.
\end{proof}

Corpus item [153] splits a unit factorial into complete residue blocks
and a remainder, then uses an exact block identity to propagate
divisibility through successive prime powers. The fiber question
requires an additional digit of precision because the leading residue of each
increment must be a unit. Below it is computed directly for the
factorial sum. The dyadic computation also requires cancellation in a
second difference.

\subsection{Odd primes}
Put $B_a(n)=S_p(n+p^a)-S_p(n)$ for $a\ge1$.
The quotient $(pk)!/(p^kk!)$ is the product of all integers up to $pk$
coprime to $p$. Passing from $k$ to $k+p^a$ appends every unit modulo
$p^{a+1}$ exactly once. Their product is $-1$ by inverse pairing, since the only self-inverse units are $1$ and $-1$. The sign in $u_p$
changes by $(-1)^{p^{a-1}}=-1$. Hence
\begin{equation}\label{fac:unit}
 u_p(k+p^a)\equiv u_p(k)\pmod{p^{a+1}}.
\end{equation}
Wilson's congruence gives $u_p(k)\equiv(-1)^{k\bmod p}\pmod p$,
so $B_1(n)\equiv\sum_{s=0}^{p-1}(-1)^s=1\pmod p$ for every $n$.
Decomposing a block into $p$ subblocks and applying \eqref{fac:unit}
term by term gives
\[
 B_{a+1}(n)\equiv p B_a(n)\pmod{p^{a+1}},\qquad
 B_a(n)\equiv p^{a-1}\pmod{p^a}.
\]
Lemma~\ref{fac:fiber}, with $h=1$, $d=0$ and $b=1$, proves the
fiber assertion for the continuous extension of $S_p$.

\subsection{The dyadic cancellation}
Write $u(k)=u_2(k)=(-1)^{\lfloor k/2\rfloor}(2k-1)!!$, taking
$u(0)=1$. For $a\ge5$ set $L=2^a$, and in this subsection let
$B_a(n)=S_2(n+L)-S_2(n)$.
\begin{lemma}\label{fac:block}
For $a\ge5$ and $n\ge0$ one has
\[
 B_a(n+L)-B_a(n)\equiv0\pmod{32L},\qquad
 B_{a+1}(n)\equiv2B_a(n)\pmod{32L}.
\]
\end{lemma}
\begin{proof}
The odd integers from $1$ to $2L-1$ have product $1$ modulo $2L$.
Inverse pairing leaves the four units $1,-1,1+L,-1+L$, whose product
is $1$ modulo $2L$. These exhaust the involutions since among two
consecutive even integers $z-1,z+1$ exactly one has valuation $1$.
As $4\mid L$, the sign of $u(L)$ is positive. Define the integer
$c=(u(L)-1)/(2L)$. The exact quotient
\[
 \frac{u(k+L)}{u(k)}=u(L)\prod_{i=1}^k
 \left(1+\frac{2L}{2i-1}\right)
\]
then gives
\begin{equation}\label{fac:difference}
 \frac{u(k+L)-u(k)}{2L}
 \equiv u(k)\left(c+\sum_{i=1}^k(2i-1)^{-1}\right)\pmod{16}.
\end{equation}
Indeed $(2L)^2$ is divisible by $32L$. The right-hand side is
$16$-periodic in $k$. In sixteen successive factors the odd residues
modulo $16$ appear twice, with product $1$ and inverse sum $0$
modulo $16$, and the sign has period four. Put $H(k)=\sum_{i=1}^k(2i-1)^{-1}$.
For $k=0,\ldots,7$ the residues modulo $16$ are
\[
\begin{array}{c|rrrrrrrr}
 k &0&1&2&3&4&5&6&7\\ \hline
 u(k)&1&1&13&1&9&1&5&1\\
 H(k)&0&1&12&9&0&9&12&1
\end{array}
\]
Both rows repeat at $k+8$. Eight successive odd residues run through
all units modulo 16, with product 1 and inverse sum 0.
Consequently the sixteen-term sums of $u$ and $uH$ are respectively
$0$ and $8$ modulo $16$. Every interval of length $L$ contains an
even number of these sixteen-term periods. Summing
\eqref{fac:difference} over $k=n+1,\ldots,n+L$ proves the first
congruence. The identity $B_{a+1}(n)=B_a(n)+B_a(n+L)$ gives the
second.
\end{proof}

\begin{proof}[Completion of the dyadic count]
At $a=5$ the first congruence in Lemma~\ref{fac:block} makes
$B_5(n)$ periodic modulo $512$ with period $32$.
We need the following initial congruences:
\begin{align*}
 &B_5(r)\equiv0\pmod{64} &&(0\le r<32),\\
 &S_2(r)\not\equiv0\pmod{64} &&(0\le r<32,\ r\notin\{0,16\}),\\
 &S_2(r)\equiv0\pmod{256},\quad B_5(r)\equiv256\pmod{512}
       &&(r=0,16).
\end{align*}
They follow by successively multiplying the unsigned factorial by
$2k-1$. The final two lines of this exact verifier also record the
valuation on the thirty excluded classes.

\begin{verbatim}
partial = [0]
odd_product = 1
for k in range(1, 64):
    odd_product = (odd_product * (2*k-1)) % 512
    term = (-1)**(k//2) * odd_product
    partial.append((partial[-1] + term) % 512)
for r in range(32):
    block = (partial[r+32] - partial[r]) % 512
    assert block % 64 == 0
    if r in (0, 16):
        assert partial[r] % 256 == 0 and block == 256
    else:
        assert partial[r] % 64 != 0
        s = (r & -r).bit_length() - 1
        v = 2*s - min(1, s)
        assert partial[r] % (1 << (v+1)) == 1 << v
\end{verbatim}

By induction the second congruence in Lemma~\ref{fac:block} gives
$B_a(n)\equiv2^{a-5}B_5(n)\pmod{2^{a+4}}$. In particular
$2^{a+1}\mid B_a(n)$ everywhere. Telescoping such increments shows
uniform continuity of $S_2$ and hence the existence of its unique
continuous extension. On each of the thirty excluded balls, all
values have the same nonzero residue modulo $64$ as their initial
value. On the two remaining balls we have
$B_{5+j}(n)\equiv2^{8+j}\pmod{2^{9+j}}$.
Lemma~\ref{fac:fiber} with $h=5,d=8,b=1$ maps each such ball
bijectively onto $256\mathbb Z_2$. Exactly one point in each maps
to zero, and $S_2(0)=0$.
\end{proof}

\subsection{Identification of the series and an exact valuation}
For $\Delta f(n)=f(n+1)-f(n)$ the identity $\Delta S_p(n)=u_p(n+1)$
implies $\Delta^m S_p(0)=\Delta^{m-1}u_p(1)$ for $m\ge1$.
Expanding the difference gives the inner coefficient of
\eqref{fac:mahler}, including its sign. Mahler's interpolation theorem \cite{Mahler} applies to the continuous function constructed above, yields
uniform convergence, and identifies its sum with that function.
This completes the proof of Theorem~\ref{fac:main}.

\begin{corollary}\label{fac:twoformula}
Let $\alpha$ be the unique nonzero dyadic root. Then for all
$x\in\mathbb Z_2$,
\[
 v_2(F_2(x))=v_2(x)+v_2(x-\alpha)-\min\{1,v_2(x)\}.
\]
\end{corollary}
\begin{proof}
On $0+32\mathbb Z_2$ and $16+32\mathbb Z_2$, the distance identity
of Lemma~\ref{fac:fiber} gives the valuation as
$3+v_2(x)$ and $3+v_2(x-\alpha)$ respectively. In each case the
distance to the other root has valuation four. On the other balls,
$v_2(x)=v_2(x-\alpha)=s\le3$. The finite recurrence above gives
$v_2(S_2(r))=2s-\min(1,s)$ for their representatives $1\le r<32$.
These valuations are at most five and hold throughout each ball by
the uniform congruence modulo $64$. At each root we use
$v_2(0)=+\infty$.
\end{proof}

\paragraph{The coefficient question.}
Keith Conrad pointed out in his answer to \cite{John} that the last unit Mahler coefficient alone cannot count the roots. In the follow-up \cite{John2}, the coefficients differ from $\Delta^m S_p(0)$ by a global sign. John C proved the initial unit range and the subsequent divisibility by $p$, then supplied an argument for a stronger odd-prime coefficient bound. The proof above obtains the leading digit of increments of the partial sum at all prime-power scales, including the dyadic cancellation. These inputs determine individual residue-class fibers.
\section{Provenance and research criteria}
\label{sec:status}
The primary target is Akbari et al., Conjecture~9. The following assessment refers to that target. The cited version is arXiv:2402.12884v2, dated 20 August 2024. Its first version appeared in February 2024.
\begin{center}
\small\renewcommand{\arraystretch}{1.23}
\begin{tabularx}{\linewidth}{@{}lX@{}}
\toprule
Criterion & Finding \\
\midrule
G1. Prior formulation & Section~5 of \cite{AGGHT} labels the universal bound Conjecture~9. The next paragraph proposes $c=1$ and raises the optimal-coefficient question. \\
G2. Same problem & Taking the indicator of an ordinary maximum matching leaves $G$, $R(G)$, $k=\nu(G)$ and the original quantifiers intact. \\
G3. Open status & The source proves the nearly-perfect case and poses the general case. The related upper-bound theorem \cite{Liu} and the literature located in the search do not settle its conjecture. Historical status is supported by the sources examined. \\
G4. Unconditional proof & Theorems~\ref{graph:main} and~\ref{graph:limit} cover all finite simple graphs and the full limit. The only external graph inequality needed by the finite proof is the classical small-mass order bound of Bollob\'as and Erd\H{o}s. \\
G5. Substantive extension & The two parts of \eqref{graph:bound} control fractional matchings at arbitrary unmatched-vertex density. The complementary thresholds in the limiting argument recover the constant in the windmill construction. Near equality yields an inverse structural conclusion. \\
G6. Priority & The literature review located no earlier proof of the general sharp bound, limiting coefficient or stability statement. The historical originality of the threshold argument and worldwide priority remain unverified. \\
\bottomrule
\end{tabularx}
\end{center}
\paragraph{Scope of the review.}
We used the sources current in the review of 16 November 2026. Targeted searches used the original titles, authors, conjecture numbers, fractional matching, generalized windmill and its limiting coefficient, and factorial Mahler roots and their related public question numbers. For comparison within the requested corpus, the searchable text of its 190 numbered TeX files was checked at the commit in \cite{VMRC}. In particular, [127] already gives the principal-eigenfunction inflection theorem of Li, Morgan and Siegrist together with the answer about the least-modulus zero. Neither is counted here as a discovery. The publicly recorded general-forcing counterexample and the suggestion of separated pulses in \cite{Omega} precede our interpolation argument. The latter prescribes exact values at predetermined sites with a positive strictly decreasing polynomial. For the arithmetic application the antecedents are identified in Section~\ref{fac:section}. The secondary questions date from 2026.
\par
These proofs are offered for independent review. No historical-first claim for an inequality, a selection principle, or a proof technique follows from the results of a bibliographic search.
\begin{thebibliography}{10}\small
\bibitem{AGGHT} S. Akbari, S. Ghasemi Nezhad, R. Ghazizadeh, J. Haslegrave and E. Tohidi. Lower bounds for the Randi\'c index in terms of matching number. 2024, v2, Section~5, Proposition~7, Theorem~8 and Conjecture~9. \url{https://arxiv.org/abs/2402.12884v2}
\bibitem{BE} B. Bollob\'as and P. Erd\H{o}s. Graphs of extremal weights. \emph{Ars Combinatoria} \textbf{50} (1998), 225--233.
\bibitem{Liu} P. Liu, F. Nan, S. O and R. Zheng. A sharp Randi\'c bound for K\H{o}nig--Egerv\'ary graphs and a conjecture of Aouchiche, Hansen, and Zheng. 2026. \url{https://arxiv.org/abs/2607.23918}
\bibitem{John} John C. Number of $p$-adic roots for the Mahler expansion of an alternating factorial sum. MathOverflow, question 509609, 28 March 2026, with an answer by K. Conrad. \url{https://mathoverflow.net/q/509609}
\bibitem{John2} John C. Lower bound on $p$-adic valuation of binomial sum. MathOverflow, question 510730, 27 April 2026, and answer 511184, 11 May 2026. \url{https://mathoverflow.net/a/511184}
\bibitem{Mahler} K. Mahler. An interpolation series for continuous functions of a $p$-adic variable. \emph{J. Reine Angew. Math.} \textbf{199} (1958), 23--34. \url{https://doi.org/10.1515/crll.1958.199.23}
\bibitem{LMS} Y. Li, R. Morgan and M. Siegrist. Concavity structure of fractional differential equations and zeros of Mittag--Leffler functions. \emph{The Journal of Analysis}, published 13 October 2026. \url{https://doi.org/10.1007/s41478-026-01192-9}
\bibitem{Omega} The Omega Institute. Li--Morgan--Siegrist 2026 inflection points, Conjecture 6.1. Public problem record, including the counterexample, the decreasing-forcing discussion and the open item on unbounded inflection numbers. \href{https://www.trureturing.org/research/li-morgan-siegrist-2026-inflection-points-conjecture-6-1/}{Problem record and certificate}.
\bibitem{VMRC} \emph{Verifiable Mathematical Reasoning Corpus}. Entries [123], [127], [153], [183] and [189]. Snapshot \texttt{5887cdea4898f0a3aa5bf7a8e0829697bd2b0f4b}, 15 November 2026. \href{https://github.com/QEEDDD/Verifiable-Mathematical-Reasoning-Corpus/tree/5887cdea4898f0a3aa5bf7a8e0829697bd2b0f4b}{Repository at the cited commit}.
\end{thebibliography}
\end{document}
